\subsection{Dirichlet interpolation on the actual finite boundary}\label{subsec:fc-line-dirichlet}

\begin{lemma}[Separation of distinct hole boundaries]\label{lem:fc-hole-boundary-separation}
\lean{EnvelopeProperty.Berkovich.RelativeSpectrum.FCStandardAffinoid.holeBoundaryPoint, EnvelopeProperty.Berkovich.RelativeSpectrum.FCStandardAffinoid.holeBoundaryPoint_mem_boundary, EnvelopeProperty.Berkovich.RelativeSpectrum.FCStandardAffinoid.holeRadius_le_center_distance, EnvelopeProperty.Berkovich.RelativeSpectrum.FCStandardAffinoid.holeRadius_lt_center_distance_of_boundary_ne, EnvelopeProperty.Berkovich.RelativeSpectrum.FCStandardAffinoid.holeRadius_lt_outer_of_boundary_ne}
\leanok
\uses{def:fc-standard-affinoid, lem:fc-inflation-linear}
Distinct open holes have centre distance at least each of their radii.
If their Gauss boundary points are distinct and $r_i\le r_j$, then
$r_i<|a_i-a_j|$. A hole boundary point distinct from the outer boundary has
radius strictly smaller than the outer radius.
\end{lemma}

\begin{proof}
If $|a_i-a_j|<r_j$, the actual evaluation point at $a_i$ belongs to both
positive-radius open holes, contradicting disjointness. Interchanging the
holes gives the other inequality. If the first radius is at most the second
and the centre distance were at most the first, both radii and the centre
distance would coincide. Centres in one closed disc have the same Gauss
point at its radius, contradicting the distinct-boundary hypothesis.
Finally an inner radius equal to the outer radius gives the same Gauss point
as the outer boundary by the centre bound.
\end{proof}

\begin{lemma}[Independent boundary logarithmic potentials]\label{lem:fc-dirichlet-independent}
\lean{EnvelopeProperty.Berkovich.RelativeSpectrum.FCStandardAffinoid.dirichletPotential, EnvelopeProperty.Berkovich.RelativeSpectrum.FCStandardAffinoid.dirichletPotential_outer, EnvelopeProperty.Berkovich.RelativeSpectrum.FCStandardAffinoid.dirichletPotential_coefficients_eq_zero}
\leanok
\uses{lem:fc-hole-boundary-separation, lem:fc-radial-skeleton-subset}
Choose one original hole for each of finitely many distinct non-outer boundary
points, with index set $I$. The functions
\[
 p_i(x)=\log r_0-\log|T-a_i|_x\quad (i\in I)
\]
are linearly independent as real functions on the actual affinoid $U$.
Each vanishes at the outer boundary.
\end{lemma}

\begin{proof}
Suppose $\sum_i\gamma_i p_i=0$ on $U$ and some coefficient is nonzero.
Choose a nonzero coefficient with least hole radius, say $\gamma_j$.
Its radius $r_j$ is strictly smaller than $r_0$ and, by the preceding lemma,
strictly smaller than every centre distance $|a_j-a_i|$ for another nonzero
coefficient. Since the family is finite, choose $t>r_j$ smaller than $r_0$
and all those distances. Both $\zeta_{a_j,r_j}$ and $\zeta_{a_j,t}$ belong
to $U$ by the radial-segment containment theorem. At the second point and
the first point every active coordinate except $j$ has the same value
$|a_j-a_i|$, whereas the $j$th values are $t$ and $r_j$. Subtracting the two
zero sums gives $\gamma_j(\log r_j-\log t)=0$. Positive radii and strict
monotonicity of the logarithm contradict $\gamma_j\ne0$.
At the outer boundary every hole coordinate has value $r_0$, proving the
last assertion.
\end{proof}

\begin{theorem}[The actual Dirichlet solution]\label{thm:fc-dirichlet-existence}
\lean{EnvelopeProperty.Berkovich.RelativeSpectrum.FCStandardAffinoid.dirichletPotential_sum_isHarmonic, EnvelopeProperty.Berkovich.RelativeSpectrum.FCStandardAffinoid.dirichletEvaluationMap, EnvelopeProperty.Berkovich.RelativeSpectrum.FCStandardAffinoid.dirichletEvaluationMap_injective, EnvelopeProperty.Berkovich.RelativeSpectrum.FCStandardAffinoid.dirichletEvaluationMap_surjective, EnvelopeProperty.Berkovich.RelativeSpectrum.FCStandardAffinoid.exists_isHarmonic_boundary_values, EnvelopeProperty.Berkovich.RelativeSpectrum.FCStandardAffinoid.existsUnique_isHarmonic_boundary_values}
\leanok
\uses{lem:fc-dirichlet-independent, thm:fc-harmonic-weak-maximum, cor:fc-harmonic-boundary-uniqueness}
For every real function $\beta$ on the actual finite boundary $\partial U$,
there is a unique harmonic function $h$ on $U$ with $h|_{\partial U}=\beta$.
This includes no holes, zero outer radius, coincident Gauss boundary points,
and zero-length radial segments.
\end{theorem}

\begin{proof}
Let $I=\partial U\setminus\{\zeta_{a_0,r_0}\}$ as a set of actual points,
so each distinct non-outer boundary point occurs exactly once. Choose an
original hole representing each point. Consider the square real linear map
\[
 E\colon\mathbf R^I\longrightarrow\mathbf R^I,\qquad
 E(\gamma)_j=\sum_i\gamma_i p_i(\zeta_{a_j,r_j}).
\]
If $E(\gamma)=0$, its harmonic sum vanishes at every boundary point,
including the outer point. The weak boundary maximum principle, applied to
the sum and its negative, makes it zero on all of $U$. The independence
lemma gives $\gamma=0$. Thus $E$ is injective, hence surjective because its
finite-dimensional domain and codomain coincide. Solve
$E(\gamma)_j=\beta(\zeta_{a_j,r_j})-\beta(\zeta_{a_0,r_0})$ and put
$h=\beta(\zeta_{a_0,r_0})+\sum_i\gamma_i p_i$. This is a constant plus a
finite real linear combination of the original logarithmic hole coordinates,
hence harmonic, and has exactly the prescribed boundary values.
Using actual boundary points removes duplicate coordinates and zero-length
segments automatically; for the empty index set the same argument gives
the constant solution.
The difference of two solutions is harmonic and has zero boundary values.
The weak boundary maximum principle applied to that difference and its
negative makes the difference zero on $U$, proving uniqueness as an actual
function.
\end{proof}
